%% ma: L12-B15-0007
%% lop: 12
%% chuong: 5
%% bai: 15
%% dang: 12.15.5
%% loai: TN4
%% muc: TH
%% dap_an: "B"
%% nguon: "(NBV)-ĐỀ SỐ 11-MỖI NGÀY 1 ĐỀ THI 2025.pdf (D11, MỖI NGÀY 1 ĐỀ - Đề số 11), Phần I Câu 12"
%% kien_thuc: "Bài 15 (phương trình đường thẳng trong không gian, vị trí tương đối của hai đường thẳng)"
%% trang_thai: chinh-thuc
%% ghi_chu: "Đã duyệt 13/10/2026. Đáp án gốc B đúng (giải hệ bằng sympy: vô nghiệm)."
\begin{ex}
Trong không gian $Oxyz$, cho hai đường thẳng $\Delta:\dfrac{x+2}{3}=\dfrac{y-1}{2}=\dfrac{z-2}{-3}$ và $d:\begin{cases}x=2+t\\y=-3+2t\\z=1-t.\end{cases}$ Khi đó, $\Delta$ và $d$
\choice{song song với nhau}{\True chéo nhau}{cắt nhau}{trùng nhau}
\loigiai{Vectơ chỉ phương của $\Delta$ là $\vec u=(3;2;-3)$, của $d$ là $\vec u_d=(1;2;-1)$. Hai vectơ không cùng phương nên $\Delta$ và $d$ hoặc cắt nhau hoặc chéo nhau.

Phương trình tham số của $\Delta$ là $x=-2+3s$, $y=1+2s$, $z=2-3s$. Xét hệ
\[\begin{cases}-2+3s=2+t\\1+2s=-3+2t\\2-3s=1-t.\end{cases}\]
Từ phương trình thứ nhất và thứ hai ta được $s=3$, $t=5$; thay vào phương trình thứ ba: $2-9=-7\ne1-5=-4$. Hệ vô nghiệm nên $\Delta$ và $d$ chéo nhau. Chọn B.}
\end{ex}
